% ch2_theory.tex — 第二章 理论基础
\section{理论基础}

本研究的理论基础建立在消费者行为研究的经典理论框架之上，主要包括Fishbein多属性态度理论、计划行为理论、态度--行为一致性理论以及消费者专业性理论。这些理论为本研究的变量测量、假设提出和结果解释提供了坚实的理论支撑。

\subsection{Fishbein多属性态度理论}

Fishbein多属性态度理论（Multi-Attribute Attitude Theory）由Fishbein \& Ajzen（1975）\cite{fishbein1975}提出，是消费者态度研究中最具影响力的理论框架之一。该理论认为，消费者对某一对象（如产品、品牌）的总体态度（Attitude, ATT）是其对该对象各属性信念的加权综合，具体表达为：

\begin{equation}
\text{ATT} = \sum_{i} (b_i \times e_i)
\label{eq:att}
\end{equation}

其中：
\begin{itemize}
  \item $b_i$（belief strength）表示消费者对第$i$个属性信念的确信程度，反映"我相信该产品具有某属性"的强度；
  \item $e_i$（belief evaluation）表示消费者对第$i$个属性的情感评价，反映"我认为该属性是好的/坏的"；
  \item $\sum$表示对所有属性信念进行求和。
\end{itemize}

该理论的核心价值在于将抽象的"态度"构念分解为可观测、可测量的信念成分，揭示了态度形成的微观机制。Aridor et al.（2022）\cite{aridor2022}的实地实验研究验证了信念强度对消费决策的影响权重，证实了多属性态度理论在现代电商情境中的适用性。

在实际应用中，多属性态度理论具有以下特点：

（1）\textbf{属性信念的识别}：消费者对产品的认知通常涉及多个属性维度（如质量、价格、外观、功能等），每个属性对应一条信念陈述。Zhang et al.（2022）\cite{zhang2022absa}的ABSA综述指出，细粒度的属性级情感分析能够更准确地捕捉消费者对不同产品维度的评价。

（2）\textbf{信念强度的差异化}：不同信念的确信程度存在差异，强信念对态度的影响权重更大。Sun et al.（2023）\cite{sun2023}的研究表明，基于LLM的情感强度评估能够有效区分不同信念的确信程度。

（3）\textbf{情感极性的双向性}：信念评价可以是正面的（$+1$）、负面的（$-1$）或中性的（$0$），这使得态度测量能够捕捉消费者的复杂情感。Yin et al.（2023）\cite{yin2023}的研究发现，在线评论中存在系统性的积极偏向，但负面信念对态度的影响往往更为显著。

（4）\textbf{态度的可加性}：总体态度是各属性信念的线性加权和，这一假设使得态度测量具有可操作性。Mitchell et al.（2022）\cite{mitchell2022}的多领域实证研究验证了多属性态度模型在不同产品类别中的稳健性。

\subsection{计划行为理论}

计划行为理论（Theory of Planned Behavior, TPB）由Ajzen（1991）\cite{ajzen1991}提出，是预测和解释人类行为的重要理论框架。该理论认为，行为意向（Behavioral Intention）是实际行为的最直接前因变量，而行为意向受到三个因素的共同影响：

（1）\textbf{态度（Attitude toward the behavior）}：个体对执行某一行为的正面或负面评价。态度越积极，行为意向越强。

（2）\textbf{主观规范（Subjective Norm）}：个体感知到的社会压力，即重要他人是否支持其执行该行为。

（3）\textbf{知觉行为控制（Perceived Behavioral Control）}：个体对执行该行为难易程度的感知，反映其对资源和机会的评估。

在本研究情境中，我们主要关注\textbf{态度对行为意向的影响}。Sheeran \& Webb（2022）\cite{sheeran2022}对十年间研究的元分析表明，态度与行为意向之间存在中等强度的正向关系（$r \approx 0.50$），这一关系在不同文化背景和行为类型中均得到验证。

计划行为理论在在线评论研究中的应用具有以下特点：

（1）\textbf{行为意向的多样性}：在电商情境中，行为意向包括再次购买意向（repurchase intention）和推荐意向（recommendation intention）。Lim et al.（2022）\cite{lim2022}的研究综述指出，这两类意向虽然相关，但受不同因素的调节。

（2）\textbf{态度--意向关系的情境依赖性}：Sheeran \& Webb（2022）\cite{sheeran2022}的元分析发现，态度--行为意向关系的强度受到个体差异（如专业性）和情境因素（如产品类别）的调节。

（3）\textbf{意向--行为差距}：Rocklage et al.（2023）\cite{rocklage2023}的研究指出，行为意向并非总能转化为实际行为，这一差距受到态度可及性（attitude accessibility）和态度强度（attitude strength）的影响。

\subsection{态度--行为一致性理论}

态度--行为一致性（Attitude-Behavior Consistency）是社会心理学的核心议题之一，探讨个体的态度如何转化为外显行为。Rocklage et al.（2023）\cite{rocklage2023}的研究系统梳理了影响态度--行为一致性的关键因素：

（1）\textbf{态度可及性（Attitude Accessibility）}：态度在记忆中的激活速度。可及性高的态度更容易被提取，因而对行为的预测力更强。个体通过直接体验形成的态度具有更高的可及性。

（2）\textbf{态度强度（Attitude Strength）}：态度的确定性和稳定性。强态度不易受外部信息干扰，与行为的一致性更高。

（3）\textbf{态度清晰度（Attitude Clarity）}：个体对自身态度的明确程度。态度越清晰，其对行为的指导作用越强。

（4）\textbf{直接经验的作用}：通过直接体验（如实际使用产品）形成的态度，相较于通过间接信息（如广告）形成的态度，具有更高的可及性和清晰度，因而与行为的一致性显著更强。

在在线评论情境中，评论者已经完成了产品购买和使用，其态度是基于直接体验形成的，因此具有较高的可及性和清晰度。Mitchell et al.（2022）\cite{mitchell2022}的多领域实证研究发现，消费者专业性显著调节态度--行为一致性：高专业性消费者的态度更为稳定，与行为的对应关系更为紧密。

\subsection{消费者专业性理论}

消费者专业性（Consumer Expertise）是指消费者在特定产品类别中通过实践积累的主观能力和经验。Mitchell et al.（2022）\cite{mitchell2022}在多领域实证分析中将消费者专业性定义为五个维度：

（1）\textbf{认知努力（Cognitive Effort）}：消费者在信息搜寻和决策过程中投入的认知资源。高专业性消费者能够更高效地处理产品信息。

（2）\textbf{认知结构（Cognitive Structure）}：消费者对产品知识的组织方式。高专业性消费者拥有更精细的产品知识图式（schema），能够快速识别产品属性。

（3）\textbf{分析能力（Analysis）}：消费者对产品信息的分析和评估能力。高专业性消费者能够更准确地判断产品质量。

（4）\textbf{精细化能力（Elaboration）}：消费者对产品信息的深度加工能力。高专业性消费者能够识别产品的细微差异。

（5）\textbf{记忆（Memory）}：消费者对产品相关信息的记忆能力。高专业性消费者能够回忆更多的产品细节。

其中，\textbf{品类熟悉度（Category Familiarity）}是消费者专业性最基础的维度，指消费者在特定品类中的经验积累程度。品类熟悉度通常通过消费者在该品类中的购买次数、使用时长或评论数量来衡量。

消费者专业性对态度和行为的影响体现在以下方面：

（1）\textbf{信息搜寻行为}：高专业性消费者更倾向于寻求属性层面的详细信息，而非依赖总体评价。这与多属性态度理论的方向高度一致。

（2）\textbf{态度稳定性}：高专业性消费者的态度更为稳定（attitude certainty），不易受外部信息干扰。

（3）\textbf{态度--行为一致性}：高专业性消费者的态度与行为之间的对应关系更为紧密，因为其态度是基于丰富的直接体验形成的。

（4）\textbf{评价标准的差异化}：高专业性消费者可能对产品形成更严苛的评价标准，从而对同一产品在态度上给予更高要求。

需要注意区分"消费者专业性"（consumer expertise）与"消费者知识"（consumer knowledge）的细微差别：消费者知识偏重客观认知存量（如对产品成分、功效的了解），而消费者专业性强调通过实践积累的主观能力和经验积累（如识别产品质量、做出差异化判断的能力）。

\subsection{理论整合与本研究的理论框架}

本研究整合上述四个理论，构建了以下理论框架：

（1）\textbf{态度测量}：采用Fishbein多属性态度理论，将消费者态度分解为可测量的信念成分（$b_i \times e_i$），通过LLM从评论文本中提取信念陈述、情感极性和信念强度。

（2）\textbf{态度--行为关系}：基于计划行为理论，检验多属性态度（ATT）对行为意向（BI）的预测关系；基于态度--行为一致性理论，检验ATT对评论评分（Rating）的预测关系。

（3）\textbf{调节机制}：基于消费者专业性理论，引入评论者经验（Reviewer Experience）作为调节变量，探讨其对ATT--Rating关系的调节作用。

（4）\textbf{方法论创新}：利用大语言模型（LLM）从非结构化评论文本中自动提取结构化心理变量，为多属性态度理论在大数据情境下的应用提供方法论工具。

上述理论框架的核心逻辑是：消费者对产品的态度由多个属性信念构成（Fishbein理论），这一态度影响其行为意向（计划行为理论）和外显行为（态度--行为一致性理论），且这一关系受到消费者专业性的调节（消费者专业性理论）。
